\documentclass[book.tex]{subfiles}
\begin{document}

\chapter{Adversarial Attacks}
\label{chap:ani}
\noindent\rule{11cm}{0.4pt} \\
\textbf{Prerequisites:}  \autoref{chap:molecular_hamiltonian},   \\
\textbf{Difficulty Level:} ** \\
\noindent\rule{11cm}{0.4pt}
\vspace{5mm}

%"Neural networks (NNs) have proven extremely effective interpolators in the physical sciences.  In particular, NN interatomic potentials can be trained to replicate the mapping of atomistic structure to energy (the so-called Potential Energy Surface) of expensive quantum mechanical methods. While they reduce the computational cost by orders of magnitude, NN potentials are brittle and struggle to generalize to points outside the training data (rare events) that maybe be visited when they are deployed in simulations. 

In this chapter, we describe how uncertainty quantification methods based on ensembles of differentiable models and on customized loss functions, including mean-variance estimation enable gradient-based active learning to systematically improve machine learned potentials. Because uncertainty metrics are differentiable end-to-end with respect to model inputs, it is possible to shift input atomic positions towards regions of low energy and high uncertainty, so they can be labeled with new quantum chemical simulations \cite{wang2020active}, \cite{ang2021active}, \cite{schwalbe2021differentiable}.

Adversarial attacks on uncertainty allow exploring the phase space  more efficiently than traditional approaches and helps train accurate and generalizable surrogate functions.

%Chelation of lithium with differentiable force fields. Lots of gaps in the training data which pose challenges. First setup was to choose a random sample, run MD/DFT to get new datapoint. Use neural reactive molecular dynamics.

%Excited state potentials. Diabetic neural networks (??) learn smooth diabetic surfaces. These somehow provide an adiabietic model.

%Attempts to make chemically transferrable 

%Get uncertainty from an ensemble of models. NN-MD is very unstable since MD is very expensive.

%Make the uncertainty itself differentiable? That provides a gradient of uncertainy. Construct a partition function based on points visited so far and use that to construct a probability density. In a 2D example, random tends towards high energy points which can be avoided by the active learning approach.

%Can do uncertainty attribution for condensed phase as well.

The loss is given by
\begin{align}
\mathcal{L} &= \frac{1}{N} \sum_{i=1}^N [\alpha_E \|E_i -\hat{E}_i\|^2 + \alpha_F \|F_i - \hat{F}_i\|^2]
\end{align}
We can compute the mean and standard deviation of the energy as
\begin{align}
\bar{E}(X) &= \frac{1}{M} \sum_{m=1}^M \hat{E}^{(m)}(X) \\
\sigma_E^2(X) &= \frac{1}{M-1} \|\hat{E}^{(m)}(X) - \bar{E}(X)\|^2
\end{align}
Similar calculations can be performed for forces. New adversarial geometries are selected by perturbing the input system
\begin{align}
\min_\theta \mathbb{E}_{(X, E, F) \sim \mathcal{D}} \left [ \max_{\delta \in \Delta} \mathcal{L}(X_\delta, E_\delta, F_\delta, \theta) \right ]
\end{align}

where $\Delta$ is a set of potential perturbations usually defined as
\begin{align}
\Delta &= \{ \delta \in \mathbb{R} | \|\delta\|_p \leq \epsilon \}
\end{align}
Adversarial updates are computed by
\begin{align}
X_\delta &= (Z, R + \delta)
\end{align}
\begin{align}
\max_{\delta \in \Delta} \mathcal{L}_{adv}(X, \delta, \theta) &= \max_{\delta \in \Delta} \sigma^2_F(X_\delta)
\end{align}
We can also Boltzmann-weight to obtain
\begin{align}
\max_{\delta \in \Delta} \mathcal{L}_{adv}(X, \delta, \theta) &= \max_{\delta \in \Delta} p(X_\delta) \sigma^2_F(X_\delta)
\end{align}
This adversarial loops allows construction of perturbed geometries that have low energy by construction but high uncertainty. Reparameterizing the force field including these new examples can lead to more robust behavior.

\end{document}